% Generated draft: requires mathematical and source review.
\documentclass[11pt]{article}
\usepackage[margin=25mm]{geometry}
\usepackage{fontspec}
\setmainfont{Latin Modern Roman}
\usepackage{amsmath,amssymb,mathtools,unicode-math}
\setmathfont{Latin Modern Math}
\usepackage{enumitem,needspace,xurl,hyperref}
\setlength{\parindent}{0pt}
\setlength{\parskip}{5pt}
\setlength{\emergencystretch}{3em}
\newcommand{\N}{\mathbb{N}}
\newcommand{\Z}{\mathbb{Z}}
\newcommand{\Q}{\mathbb{Q}}
\newcommand{\R}{\mathbb{R}}
\newcommand{\C}{\mathbb{C}}
\newcommand{\F}{\mathbb{F}}
\newcommand{\A}{\mathbb{A}}
\newcommand{\PP}{\mathbb{P}}
\newcommand{\E}{\mathbb{E}}
\newcommand{\eps}{\varepsilon}
\newcommand{\dd}{\,\mathrm{d}}
\begin{document}
\section*{UnsolvedMath ID: 30005540}
\textbf{Source problem code:} catalogue entry
\subsection{Definitions and mathematical statement}
The supplied record attributes the following conjecture to Singer. Its precise hypotheses require clarification: it does not specify compactness, boundary conditions, or whether ``Yamabe invariant'' means the invariant of the conformal class or that of the underlying manifold.

To describe the usual closed-manifold interpretation explicitly, suppose that \(M\) is a smooth, closed, oriented four-manifold and \(g\) is a smooth Riemannian metric. These domain assumptions are provisional additions for this interpretation, not hypotheses explicitly supplied by the record. The Hodge star gives \(\Lambda^2T^*M=\Lambda_g^+\oplus\Lambda_g^-\). Write \(W_g^+\) for the self-dual Weyl curvature; the anti-self-duality condition is \(W_g^+=0\).

Let \(\Gamma(E)\) denote smooth sections of a bundle \(E\), and let \(S_0^2E\) denote its trace-free symmetric square. At an anti-self-dual metric, the linearized self-dual Weyl curvature operator is
\[
D_g:\Gamma(S_0^2T^*M)\longrightarrow\Gamma(S_0^2\Lambda_g^+),\qquad D_gh=\left.\frac{\dd}{\dd t}\right|_{t=0}W_{g+th}^+,
\]
using the metric identifications of the varying target bundles. Let \(D_g^*\) be its formal adjoint for the \(L^2\) inner products induced by \(g\). In the terminology used by the supplied statement, unobstructedness means that the obstruction space \(\ker D_g^*\) vanishes.

For the closed-manifold interpretation, define the conformal Yamabe constant by
\[
Y(M,[g])=\inf_{\substack{u\in C^\infty(M)\\u>0}}
\frac{\displaystyle\int_M\bigl(6|\nabla_gu|_g^2+R_gu^2\bigr)\,\dd V_g}
{\displaystyle\left(\int_Mu^4\,\dd V_g\right)^{1/2}},
\]
where \(R_g\) is the scalar curvature. The smooth Yamabe invariant of the manifold is instead
\[
\sigma(M)=\sup_{[h]}Y(M,[h]),
\]
where the supremum ranges over smooth Riemannian conformal classes on \(M\).

\textbf{Conjectural target, with unresolved interpretation.} If the source means positivity in the given conformal class, the target is
\[
\forall (M,g)\text{ as above},\qquad
W_g^+=0\ \text{and}\ Y(M,[g])>0
\quad\Longrightarrow\quad \ker D_g^*=\{0\}.
\]
If it means the manifold Yamabe invariant, the hypothesis is instead \(\sigma(M)>0\). These are distinct formulations; the supplied record does not determine which is intended. If noncompact manifolds or manifolds with boundary are intended, the Yamabe quantity, section spaces, adjoint realization, and boundary or decay conditions must also be specified. No unconditional choice among these formulations is made here.
\subsection{Short English statement}
The conjecture attributed to Singer asks whether positivity of the relevant Yamabe quantity forces an anti-self-dual four-manifold to have zero obstruction space \(\ker D_g^*\). The supplied wording leaves both the geometric domain and the meaning of Yamabe positivity unresolved. Its supplied status is partially solved; this draft does not independently verify that assessment.
\subsection{Sources}
\begin{itemize}
\item Local registry: supplied UnsolvedMath source record, permanent ID \(30005540\), source problem code \texttt{catalogue entry}, titled ``Singer's Unobstructedness Conjecture for Anti-Self-Dual Four-Manifolds.'' The record attributes the conjecture to Singer and gives proposed year 2023; no proposer is separately named.
\item Supplied report citation: \emph{Hypoelliptic Operators in Geometry} (2023), Oberwolfach Reports, pp.\ 1323--1396, DOI 10.4171/owr/2023/24. \url{https://doi.org/10.4171/owr/2023/24}. The record describes Oberwolfach Reports as documenting workshops and their open-problem discussions.
\item Supplied literature source, ``OWR 2023 24'': \url{https://publications.mfo.de/bitstream/handle/mfo/4050/OWR_2023_24.pdf?isAllowed=y&sequence=4}.
\item Supplied literature source, ``Gover Gursky'': \url{https://doi.org/10.1515/crelle-2024-0028}. The supplied assessment says that Gover and Gursky establish unobstructedness under additional conformal and topological inequalities, while the full positive Yamabe Singer conjecture remains unproved. Those inequalities are not stated in the record.
\item Literature assessment date: \texttt{2026-08-22}. The partially solved label and the preceding assessment are supplied dated curation derived from the OWR material, not independently verified here; the record itself requests independent verification of time-sensitive status and cited claims.
\item The registry labels the statement \texttt{corrected\_verified} and says it was reconstructed and checked against the official report. This is supplied provenance, not an independent verification by this draft. The original-statement field contains only an attribution to Singer and a reference to ``Conjecture 1,'' rather than the actual report passage.
\end{itemize}
\end{document}
